\documentclass[13pt, a4paper]{extarticle} % Hỗ trợ cỡ chữ 13pt chuẩn BM18

% ----- CÀI ĐẶT FONT VÀ NGÔN NGỮ (DÙNG PDFLATEX) -----
\usepackage[utf8]{inputenc} % Hỗ trợ gõ tiếng Việt
\usepackage[T5]{fontenc}    % Bảng mã tiếng Việt
\usepackage{mathptmx}       % Font Times mổ phỏng cho pdfLaTeX
\usepackage[vietnamese]{babel}

% ----- CÀI ĐẶT KHỔ GIẤY VÀ LỀ -----
\usepackage{geometry}
\geometry{
    a4paper,
    left=3cm, right=2cm, top=2.5cm, bottom=2.5cm % Lề tiêu chuẩn BM18
}

% ----- CÀI ĐẶT ĐÁNH SỐ TRANG GIỮA BÊN TRÊN -----
\usepackage{fancyhdr}
\pagestyle{fancy}
\fancyhf{} % Xóa mọi định dạng header/footer cũ
\fancyhead[C]{\thepage} % Đánh số trang ở chính giữa phía trên
\renewcommand{\headrulewidth}{0pt} % Xóa đường kẻ ngang ở header

% ----- CÁC GÓI TOÁN HỌC VÀ HÌNH ẢNH DÀNH CHO VRPTW -----
\usepackage{amsmath, amssymb}
\usepackage{graphicx}
\usepackage{caption}
\usepackage{booktabs}
\usepackage{indentfirst} % Thụt lề đoạn đầu tiên

% ----- CÀI ĐẶT TRÍCH DẪN (AUTHOR-YEAR) -----
\usepackage[authoryear, round]{natbib}
\renewcommand{\bibname}{Tài liệu tham khảo}

\begin{document}

% ================= TRANG BÌA =================
\begin{titlepage}
    \thispagestyle{empty} % TẮT ĐÁNH SỐ TRANG CHO TRANG BÌA
    \begin{center}
        TỔNG LIÊN ĐOÀN LAO ĐỘNG VIỆT NAM \\
        TRƯỜNG ĐẠI HỌC TÔN ĐỨC THẮNG \\
        --------------------------------- \\
        \vspace{2.5cm}
            ĐÁNH THỬ, COI NHƯ NHÁP, KHÔNG DÙNG CÁI NÀY ĐƯA CÔ COI, NĂN NỈ NĂN NỈ NĂN NỈ
            \vspace{1.5cm}
        
        Công trình Nghiên cứu khoa học sinh viên năm học 2025 - 2026 \\
        \vspace{1.5cm}
        
        TÊN ĐỀ TÀI \\
        \vspace{0.5cm}
        \textbf{\Large BỘ ĐIỀU KHIỂN CHỐNG CHỮNG DỰA TRÊN HỌC TĂNG CƯỜNG SÂU CHO ALNS GIẢI VRPTW (DDQN-ALNS)} \\
        \vspace{1.5cm}
        
        ĐƠN VỊ: KHOA CÔNG NGHỆ THÔNG TIN
        \vspace{2.5cm}
    \end{center}
    
    \begin{flushleft}
        GIẢNG VIÊN HƯỚNG DẪN: \\
        TS. Hồ Thị Linh \\
        \vspace{0.8cm}
        SINH VIÊN THỰC HIỆN: \\
        1. Huỳnh Nhật Huy \\
        2. Nguyễn Nhật Huy \\
        3. .......... Trân

    \end{flushleft}
    
    \vfill
    \begin{center}
        TP. Hồ Chí Minh, tháng 4 năm 2026
    \end{center}
\end{titlepage}

\newpage
\setcounter{page}{1}

%========================================================== INTRODUCTION
\section*{GIỚI THIỆU VÀ ĐẶT VẤN ĐỀ}
Trong bối cảnh logistics và thương mại điện tử, bài toán VRPTW (Vehicle Routing Problem with Time Windows) đóng vai trò trung tâm trong tối ưu chi phí vận chuyển và tuân thủ cam kết thời gian giao hàng. Adaptive Large Neighborhood Search (ALNS) là metaheuristic phổ biến nhờ khả năng phá huỷ--tái tạo linh hoạt, nhưng dễ rơi vào trạng thái \textit{plateau}: nhiều vòng lặp liên tiếp không cải thiện nghiệm tốt nhất do phân phối xác suất chọn toán tử đã hội tụ về một pha khai thác cục bộ. Hiện tượng này làm lãng phí ngân sách tính toán và cản trở việc thoát khỏi cực trị địa phương.

Nghiên cứu này đề xuất bộ giải \textit{PlateauHybridSolver} (DDQN-ALNS) trong đó một Double DQN hoạt động như \textbf{bộ điều khiển mức pha}: chỉ can thiệp khi phát hiện chững (\texttt{no\_imp} $\ge 60$) và lựa chọn một trong năm chế độ tìm kiếm tổng quát để hiệu chỉnh nhiệt độ, kích thước phá huỷ và phân phối toán tử. Mục tiêu là đạt nghiệm tốt hơn ALNS thuần tuý về tổng quãng đường (TD) \textbf{và} rút ngắn thời gian chạy, đồng thời đánh giá khả năng \textit{zero-shot transfer} của chính sách đã học.

\newpage

%========================================================== 
\section{Tổng quan tài liệu}
Phương pháp Quy hoạch tuyến tính của \citep{Dantzig1959} đặt nền móng cho VRP, còn bộ dữ liệu chuẩn Solomon \citep{Solomon1987} được dùng rộng rãi để benchmark VRPTW. ALNS được giới thiệu bởi Ropke và Pisinger làm khung phá huỷ--tái tạo thích ứng; biến thể SISR và các phát triển gần đây tập trung vào toán tử có tri thức miền. Hướng đang nổi là dùng Học tăng cường sâu điều khiển metaheuristic: \citep{Chen2020Hybrid} dùng DQN chọn toán tử ở cấp vòng lặp, còn nhiều công trình sau đó (ví dụ kết hợp GNN với PPO) mã hoá trực tiếp cấu trúc đồ thị. Khác với hướng chọn toán tử trực tiếp, nghiên cứu này đặt DDQN ở \textbf{mức pha} (mode-level) chỉ kích hoạt khi ALNS chững -- giúp giảm chi phí huấn luyện và cho phép tái sử dụng chính sách qua các bộ dữ liệu.

\section{Mục tiêu - Phương pháp}

\subsection{Mục tiêu nghiên cứu}
Ba mục tiêu đo lường được:
\begin{enumerate}
    \item Giảm sai số tương đối tổng quãng đường (Gap\%) so với Best Known Solution trên các bộ Solomon RC1 và RC2 ở mức có ý nghĩa thống kê (Wilcoxon $p < 0.05$).
    \item Rút ngắn thời gian chạy trung bình so với ALNS thuần tuý trên cùng ngân sách vòng lặp.
    \item Đánh giá khả năng zero-shot transfer: chính sách DDQN huấn luyện trên RC1 được áp dụng trực tiếp cho RC2 mà không huấn luyện lại.
\end{enumerate}

\subsection{Kiến trúc DDQN-ALNS}

Hệ thống được mô hình hoá dưới dạng Quá trình Quyết định Markov (MDP) với ba thành phần cốt lõi: trạng thái $S_t$, hành động (mode) $A_t$, và phần thưởng $R_t$. Vòng lặp ALNS được chia thành các \textit{segment} có độ dài 50 vòng; sau mỗi segment, DDQN quan sát trạng thái tổng hợp và ra quyết định mode cho segment kế tiếp. DDQN \textbf{chỉ can thiệp} khi \texttt{no\_imp} $\ge$ \texttt{plateau\_start} = 60; ngoài điểm chững, ALNS vận hành ở chế độ mặc định để giữ tốc độ khám phá.

\subsubsection{Không gian Trạng thái (10 chiều)}
Để chính sách không phụ thuộc vào kích thước đồ thị, vector trạng thái được chuẩn hoá về $[0, 1]$:
\[
S_t = \big[\text{patience},\, \text{cost\_gap},\, \text{temp\_ratio},\, \text{imp\_rate},\, \text{nv\_ratio},\, \text{route\_balance},\, \text{load\_balance},\, \text{tw\_tight},\, \text{avg\_slack},\, \text{progress}\big]
\]
trong đó \texttt{patience} là tỉ lệ \texttt{no\_imp} trên ngưỡng dừng sớm, \texttt{cost\_gap} là khoảng cách tương đối giữa nghiệm hiện tại và nghiệm tốt nhất, \texttt{temp\_ratio} là nhiệt độ chuẩn hoá, \texttt{imp\_rate} là tỉ lệ cải thiện trong 50 vòng gần nhất, \texttt{route\_balance} và \texttt{load\_balance} mô tả độ đồng đều của các tuyến, \texttt{tw\_tight} và \texttt{avg\_slack} nắm bắt độ ngặt nghèo của ràng buộc thời gian.

\subsubsection{Không gian Hành động (5 mode)}
Thay vì chọn trực tiếp toán tử cơ sở, DDQN đưa ra một trong năm chế độ tìm kiếm:
\begin{itemize}
    \item \textbf{Default}: phân phối xác suất ALNS gốc.
    \item \textbf{Intensify}: thu nhỏ kích thước phá huỷ (scale $0.70$), ưu tiên toán tử \textit{worst}/\textit{shaw removal} để khai thác sâu vùng hiện tại.
    \item \textbf{Diversify}: tăng kích thước phá huỷ (scale $1.20$) và tăng nhiệt độ (boost $1.08$) để thoát cực trị.
    \item \textbf{TW Rescue}: tăng trọng số cho \textit{time-oriented removal} nhằm sửa vi phạm cửa sổ thời gian.
    \item \textbf{Route Reduce}: ưu tiên \textit{shortest-route removal} và \textit{regret-k insertion} nhằm ép tuyến yếu nhất biến mất khi độ đầy đội xe thấp.
\end{itemize}
Mode \textit{Route Reduce} được gọi tự động (không qua DDQN) khi phát hiện plateau kèm theo độ đầy đội xe dưới ngưỡng động $\max(0.52,\, 0.80 - 0.25 \cdot \text{tw\_tight})$.

\subsubsection{Hàm Phần thưởng theo segment}
Gọi $\Delta_{\text{best}}^{\text{NV}} = \text{NV}_{\text{best,before}} - \text{NV}_{\text{best,after}} \ge 0$ khi nghiệm tốt nhất giảm được số xe, và $r_{\text{best}}^{\text{TD}} = (\text{TD}_{\text{best,before}} - \text{TD}_{\text{best,after}})/\text{TD}_{\text{best,before}}$ là mức cải thiện tương đối. Tương tự cho nghiệm hiện tại $\text{cur}$. Phần thưởng mỗi segment:
\[
R_t = -0.75 \;-\; 0.75 \cdot \mathbb{I}[\text{accepted}=0] \;+\; R^{\text{best}} \;+\; R^{\text{cur}} \;+\; 3\big(\rho_{\text{after}} - \rho_{\text{before}}\big)
\]
với $\rho$ là độ đầy đội xe (fleet fill), và
\[
R^{\text{best}} = \begin{cases}
30\,\Delta_{\text{best}}^{\text{NV}} + 100\max(0,\, r_{\text{best}}^{\text{TD}}) & \Delta_{\text{best}}^{\text{NV}} > 0 \\
400\, r_{\text{best}}^{\text{TD}} & \Delta_{\text{best}}^{\text{NV}} = 0,\; r_{\text{best}}^{\text{TD}} > 0 \\
0 & \text{ngược lại}
\end{cases},\quad
R^{\text{cur}} = \begin{cases}
75\, r_{\text{cur}}^{\text{TD}} & \Delta_{\text{cur}}^{\text{NV}} = 0,\; r_{\text{cur}}^{\text{TD}} > 0 \\
-20\, |\Delta_{\text{cur}}^{\text{NV}}| & \text{cur tăng xe} \\
0 & \text{ngược lại}
\end{cases}
\]
Thiết kế này đặt trọng số cao nhất cho việc cải thiện nghiệm \textbf{tốt nhất} ($R^{\text{best}}$ có hệ số gấp 4--5 lần $R^{\text{cur}}$), đồng thời phạt mức vừa phải ($-20$) khi nghiệm hiện tại thêm xe -- nhằm giữ động lực khám phá nhưng tránh trôi dạt.

\subsection{Huấn luyện}
QNet 2 lớp ẩn 96 đơn vị, Adam $\text{lr}=3\times 10^{-4}$, $\gamma=0.95$, $\varepsilon$-greedy giảm từ $0.35$ xuống $0.05$. Replay buffer $4096$ transitions, cập nhật mạng đích mỗi 20 bước. Mỗi bài toán chạy 2000 vòng ALNS (chia thành 40 segment); cấu hình đầy đủ nằm trong \texttt{Config} của mã nguồn.

\section{Kết quả - Thảo luận}

Toàn bộ kết quả được báo cáo trên 5 lần chạy độc lập với các seed khác nhau cho mỗi cặp (instance, thuật toán). BKS dùng để so sánh lấy từ thư viện Solomon chuẩn.

\begin{table}[h]
\centering
\caption{Kết quả tổng hợp trên Solomon RC1 và RC2 (trung bình 5 runs).}
\label{tab:summary}
\begin{tabular}{llrrrrr}
\toprule
\textbf{Bộ} & \textbf{Thuật toán} & \textbf{NV} & \textbf{TD} & \textbf{Gap\%} & \textbf{OnTime} & \textbf{Thời gian} \\
\midrule
RC1 & ALNS        & 12.5 $\pm$ 0.3 & 1393.5 $\pm$ 13.1 & +1.0\% & 100.0\% & 120.7s \\
RC1 & DDQN-ALNS   & 12.7 $\pm$ 0.3 & \textbf{1387.3} $\pm$ 18.4 & \textbf{+0.5\%} & 100.0\% & \textbf{79.8s} \\
\midrule
RC2 & ALNS        & 3.5 $\pm$ 0.1  & 1132.2 $\pm$ 23.8 & +1.6\% & 100.0\% & 91.3s \\
RC2 & DDQN-ALNS   & 3.6 $\pm$ 0.1  & 1121.0 $\pm$ 32.3 & \textbf{+0.7\%} & 100.0\% & 64.4s \\
RC2 & DDQN-ALNS$^\star$ & 3.6 $\pm$ 0.1 & 1121.4 $\pm$ 31.9 & +0.7\% & 100.0\% & 64.2s \\
\bottomrule
\end{tabular}
\vspace{0.2em}

\small $\star$ = chính sách transfer huấn luyện trên RC1, áp dụng zero-shot cho RC2.
\end{table}

\subsection{Chất lượng nghiệm (Gap\%)}
Trên RC1, DDQN-ALNS giảm Gap\% trung bình từ $+1.0\%$ xuống $+0.5\%$. Kiểm định Wilcoxon signed-rank (hai phía) trên 8 instance cho $W=0.0,\, p=0.0078$ -- \textbf{có ý nghĩa thống kê ở mức} $\alpha = 0.05$. Trên RC2 xu hướng tương tự (Gap từ $+1.6\%$ xuống $+0.7\%$) nhưng không đạt ngưỡng ($p=0.1953$) do cỡ mẫu 8 instance nhỏ và phương sai cao hơn. Về số xe NV, hai thuật toán gần như tương đương (chênh lệch trung bình $\le 0.2$ xe, không có ý nghĩa thống kê): DDQN-ALNS không giảm thêm xe so với ALNS nhưng cũng không \textit{đánh đổi} xe để giảm quãng đường -- phù hợp với thiết kế hàm thưởng ưu tiên nghiệm tốt nhất và phạt tăng xe.

\subsection{Thời gian chạy}
Đóng góp thực tiễn rõ rệt nhất là thời gian: RC1 giảm $120.7 \to 79.8$s ($-33.9\%$), RC2 giảm $91.3 \to 64.4$s ($-29.5\%$). Lý do: khi DDQN kích hoạt chế độ \textit{Intensify} hoặc \textit{Route Reduce} đúng lúc chững, ngân sách tính toán được tập trung cho các pha hội tụ nhanh, sớm thoả điều kiện dừng sớm (\texttt{early\_stop\_patience} = 400).

\subsection{Zero-shot transfer}
Cột DDQN-ALNS$^\star$ dùng nguyên bộ trọng số QNet huấn luyện trên RC1 để giải RC2 mà \textbf{không cập nhật gradient} nào. Kết quả gần như trùng khớp với DDQN-ALNS huấn luyện trực tiếp trên RC2 (TD $1121.4$ vs $1121.0$, Gap $0.7\%$ vs $0.7\%$), chứng minh chính sách ở mức pha có tính tổng quát hoá: đặc trưng trạng thái không phụ thuộc hình học cụ thể, do đó có thể tái sử dụng qua các cấu trúc đồ thị khác nhau.

\subsection{Hạn chế}
(i) Phân tích chỉ giới hạn trên hai bộ RC1/RC2; cần mở rộng sang C, R, và Gehring-Homberger lớn hơn. (ii) Số xe NV không được cải thiện so với ALNS; hướng mở rộng có thể thay đổi hàm thưởng hoặc thêm pha tái tổ hợp dựa trên \textit{route pool} tích cực hơn. (iii) Phân phối bài toán mục tiêu khi triển khai thực tế (dữ liệu giao hàng đô thị) khác đáng kể so với Solomon tổng hợp; cần fine-tune.

\section{Kết luận - Đề nghị}
Nghiên cứu đã trình bày DDQN-ALNS dưới dạng bộ điều khiển mức pha cho metaheuristic giải VRPTW. Trên Solomon RC1, kết quả giảm Gap\% có ý nghĩa thống kê ($p = 0.0078$) kèm rút ngắn thời gian chạy $\approx 30\%$. Chính sách đã học có tính transfer tốt sang RC2 mà không cần huấn luyện lại. Hướng tiếp theo: tích hợp Mạng Nơ-ron Đồ thị (GNN) để mã hoá đặc trưng hình học trực tiếp, mở rộng thử nghiệm sang tập RC lớn hơn, và kết hợp giảng viên (dạy--trò) giữa ALNS mạnh (Gurobi-ALNS) và DDQN nhằm giảm NV.

% ================= TÀI LIỆU THAM KHẢO =================
\newpage
\begin{thebibliography}{}

\bibitem[Dantzig và Ramser (1959)]{Dantzig1959}
Dantzig, G. B. và Ramser, J. H. (1959). \textit{The Truck Dispatching Problem}. Management Science, 6(1), 80--91.

\bibitem[Solomon (1987)]{Solomon1987}
Solomon, M. M. (1987). \textit{Algorithms for the Vehicle Routing and Scheduling Problems with Time Window Constraints}. Operations Research, 35(2), 254--265.

\bibitem[Chen và cộng sự (2020)]{Chen2020Hybrid}
Chen, X., Ulmer, M. W. và Thomas, B. W. (2020). \textit{Deep Q-Learning for Same-Day Delivery with Vehicles and Drones}. European Journal of Operational Research.

\end{thebibliography}

\end{document}
